\section*{Chapter \ref{ch:b3:galois} --- Field Extensions and
Galois Theory}

\begin{solution}{exo:b3:galois:1}
$\sqrt3 \notin \Q(\sqrt2)$: from $\sqrt3 = a + b\sqrt2$ ($a, b
\in \Q$), squaring gives $3 = a^2 + 2b^2 + 2ab\sqrt2$, so $ab =
0$; $b = 0$ makes $\sqrt3$ rational, $a = 0$ gives $\sqrt6 = 2b
\in \Q$ --- both false (standard prime-factorization arguments).
Hence $[\Q(\sqrt2,\sqrt3) : \Q(\sqrt2)] = 2$ and the tower law
gives degree $4$.

Let $\gamma = \sqrt2 + \sqrt3$. Then $\gamma^2 = 5 + 2\sqrt6$ and
$(\gamma^2 - 5)^2 = 24$: $\gamma$ annihilates $X^4 - 10X^2 + 1$.
Moreover $\gamma^3 = 11\sqrt2 + 9\sqrt3$, so $\gamma^3 - 9\gamma
= 2\sqrt2$: $\sqrt2 \in \Q(\gamma)$, then $\sqrt3 = \gamma -
\sqrt2 \in \Q(\gamma)$: $\Q(\gamma) = \Q(\sqrt2,\sqrt3)$, of
degree $4$. The annihilating quartic, having the degree of the
minimal polynomial, \emph{is} the minimal polynomial: $X^4 -
10X^2 + 1$ (in particular it is irreducible over $\Q$).
\end{solution}

\begin{solution}{exo:b3:galois:2}
The three roots of $X^3 - 2$ in $\C$ are $\alpha, j\alpha,
j^2\alpha$ with $j = \eu^{2\iu\pi/3}$. The map $\alpha \mapsto
j\alpha$ defines a $\Q$-embedding $\Q(\alpha) \to \C$
(\cref{thm:b3:galois:simple}: both generate degree-$3$ extensions
with the same minimal polynomial), whose image $\Q(j\alpha)
\not\subseteq \R$ differs from $\Q(\alpha) \subseteq \R$:
$\Q(\alpha)/\Q$ is not normal. The splitting field is $L =
\Q(\alpha, j)$, with $[L:\Q] = [L:\Q(\alpha)]\,[\Q(\alpha):\Q] =
2 \cdot 3 = 6$ ($j$ satisfies $X^2 + X + 1$, irreducible over
the real field $\Q(\alpha)$). An automorphism of $\Q(\alpha)$
must send $\alpha$ to a root of $X^3 - 2$ \emph{inside}
$\Q(\alpha) \subseteq \R$: only $\alpha$ qualifies, so
$\operatorname{Aut}_\Q(\Q(\alpha)) = \{\mathrm{id}\}$, of order
$1 < 3$.
\end{solution}

\begin{solution}{exo:b3:galois:3}
$X^2 + 1$ has no root in $\mathbb F_3$ ($0, 1, 2 \mapsto 1, 2,
2$), so $\mathbb F_9 = \mathbb F_3[X]/(X^2+1)$ is a field with
$9$ elements; write $\omega = \bar X$, $\omega^2 = -1$. The
group $\mathbb F_9^\times$ is cyclic of order $8$; $\omega$ has
order $4$, but $1 + \omega$ works: $(1+\omega)^2 = 1 + 2\omega +
\omega^2 = 2\omega$, $(1+\omega)^4 = 4\omega^2 = \omega^2 = -1
\ne 1$: order $8$.

Over $\mathbb F_2$ --- degree $1$: $X$, $X + 1$; degree $2$:
$X^2 + X + 1$ (the other three quadratics have roots); degree
$3$: $X^3 + X + 1$ and $X^3 + X^2 + 1$ (no roots in $\mathbb
F_2$; the other six cubics have roots). Verification:
\[
(X^3{+}X{+}1)(X^3{+}X^2{+}1) = X^6 + X^5 + X^4 + X^3 + X^2 + X
+ 1,
\]
and $X(X{+}1)(X^6 + \dots + 1) = X(X^7 + 1) = X^8 + X = X^8 -
X$ over $\mathbb F_2$ --- exactly the irreducibles of degree
dividing $3$, as \cref{exo:b3:galois:6} predicts (degree $2$ is
absent: $2 \nmid 3$).
\end{solution}

\begin{solution}{exo:b3:galois:4}
(a) Mod $7$: the powers of $3$ are $3, 2, 6, 4, 5, 1$: order
$6$, a generator; the primitive roots are the $3^k$ with
$\gcd(k, 6) = 1$: $3$ and $3^5 = 5$. Mod $11$: powers of $2$:
$2, 4, 8, 5, 10, 9, 7, 3, 6, 1$: a generator; primitive roots
$2^k$, $\gcd(k, 10) = 1$: $2, 2^3 = 8, 2^7 = 7, 2^9 = 6$.

(b) Write $x = g^k$ with $g$ a generator
(\cref{thm:b3:galois:cyclic}). Then $x$ is a square iff $k$ is
even (squares are the $g^{2l}$, and $g^{2l} = g^{k}$ iff $k
\equiv 2l \bmod p-1$, solvable iff $k$ even, $p - 1$ being
even). And $x^{(p-1)/2} = g^{k(p-1)/2} = 1$ iff $(p-1) \mid
k\frac{p-1}2$ iff $k$ even: the two conditions agree. For $x =
-1 = g^{(p-1)/2}$: it is a square iff $\frac{p-1}2$ is even, iff
$p \equiv 1 \pmod 4$ --- \cref{pb:b3:rings:1} again.
\end{solution}

\begin{solution}{exo:b3:galois:5}
Inside $\bar{\mathbb F}_p$, $\mathbb F_{p^k} = \{x : x^{p^k} =
x\}$ is the fixed set of $F^k$. The intersection $\mathbb F_{p^m}
\cap \mathbb F_{p^n}$ is fixed by $F^m$ and $F^n$, hence by
$F^{\gcd(m,n)}$ ($\gcd = am + bn$: on a fixed element, $F^{am +
bn} = (F^m)^a(F^n)^b$ acts trivially --- exponents may be taken
positive by periodicity); so it lies in $\mathbb
F_{p^{\gcd(m,n)}}$, which is conversely contained in both
(\cref{thm:b3:galois:finitefields}(3)). The composite $\mathbb
F_{p^m}\mathbb F_{p^n}$: any field containing both has degree
divisible by $m$ and $n$, hence by $\operatorname{lcm}(m,n)$; and
$\mathbb F_{p^{\operatorname{lcm}}}$ contains both: it is the
composite. For $m \mid n$: $\operatorname{Gal}(\mathbb
F_{p^n}/\mathbb F_{p^m})$ consists of the powers of $F$ fixing
$\mathbb F_{p^m}$, i.e.\ of $F^m$: cyclic of order $n/m$,
generated by $F^m \colon x \mapsto x^{p^m}$ (order as in
\cref{thm:b3:galois:finitefields}(2)).
\end{solution}

\begin{solution}{exo:b3:galois:6}
$X^{q^n} - X$ is separable (derivative $-1$) with root set
$\mathbb F_{q^n}$. Let $P$ be monic irreducible of degree $d$.
If $d \mid n$: $\mathbb F_q[X]/(P) \cong \mathbb F_{q^d}
\subseteq \mathbb F_{q^n}$, so $P$ has a root $\alpha \in
\mathbb F_{q^n}$; $\alpha^{q^n} = \alpha$, and $P =
\pi_\alpha$ divides $X^{q^n} - X$. If $P \mid X^{q^n} - X$: a
root $\alpha \in \mathbb F_{q^n}$ generates $\mathbb F_{q^d}
\subseteq \mathbb F_{q^n}$, so $d \mid n$
(\cref{thm:b3:galois:finitefields}(3)). Distinct irreducibles
are coprime and the product is separable: each $P$ appears with
exponent exactly $1$, and every root of $X^{q^n}-X$ is a root
of its minimal polynomial: the factorization holds. Comparing
degrees: $q^n = \sum_{d\mid n} d\,I_d(q)$.

Consequently $I_1 = q$; $q^2 = I_1 + 2I_2$ gives $I_2 =
\frac{q^2 - q}2$; $q^3 = I_1 + 3I_3$ gives $I_3 = \frac{q^3 -
q}3$; $q^4 = I_1 + 2I_2 + 4I_4$ gives $I_4 = \frac{q^4 -
q^2}4$. Existence: $nI_n = q^n - \sum_{d \mid n,\, d < n}
dI_d \geq q^n - \sum_{d \leq n/2} q^d > q^n - q^{n/2 + 1} \geq
0$ for $n \geq 2$ (and $I_1 = q \geq 1$): $I_n \geq 1$ always.
\end{solution}

\begin{solution}{exo:b3:galois:7}
$L = \Q(\sqrt2, \sqrt3)$ is the splitting field of $(X^2 -
2)(X^2 - 3)$, separable: Galois of degree $4$
(\cref{exo:b3:galois:1}). An automorphism sends $\sqrt2 \mapsto
\pm\sqrt2$ and $\sqrt3 \mapsto \pm\sqrt3$: at most $4$ choices,
and $\abs G = 4$ realizes all: $G \cong (\Z/2\Z)^2$, with
elements $\mathrm{id}, \sigma (\sqrt2 \mapsto -\sqrt2), \tau
(\sqrt3\mapsto-\sqrt3), \sigma\tau$. Subgroups of order $2$:
$\langle\sigma\rangle, \langle\tau\rangle,
\langle\sigma\tau\rangle$, with fixed fields $\Q(\sqrt3)$,
$\Q(\sqrt2)$, $\Q(\sqrt6)$ (note $\sigma\tau$ fixes $\sqrt6 =
\sqrt2\sqrt3$). The lattice: $\Q$ below, the three quadratic
fields in the middle, $L$ on top --- and nothing else
(\cref{thm:b3:galois:fundamental}). For $(X^2-2)(X^2-3)(X^2-6)$:
the splitting field is the \emph{same} $L$ ($\sqrt6 =
\sqrt2\sqrt3$), so the answer is identical: the Galois
correspondence is an invariant of the \emph{extension}, not of
the polynomial chosen to present it.
\end{solution}

\begin{solution}{exo:b3:galois:8}
(a) $G$ acts faithfully and transitively (irreducibility) on the
three roots: $G \hookrightarrow S_3$ with $3 \mid \abs G$: $G
\cong A_3$ or $S_3$. Every $\sigma \in G$ permutes the $x_i$, and
$\sigma(\delta) = \varepsilon(\sigma)\,\delta$ ($\delta$ is
alternating in the roots). If $\Delta$ is a square in $\Q$:
$\delta \in \Q^\times$ (note $\delta \neq 0$: separable), so
$\varepsilon(\sigma) = 1$ for all $\sigma$: $G \subseteq A_3$,
hence $= A_3$. If not: $\delta \notin \Q$, so some $\sigma$ has
$\varepsilon(\sigma) = -1$: $G = S_3$. (In both cases $\Q(\delta)
= L^{G \cap A_3}$.)

(b) With $x_1 + x_2 + x_3 = 0$: $u + v = 2x_1 - (x_2 + x_3) =
3x_1$. Also
\[
uv = \sum_i x_i^2 + (j + j^2)\sum_{i<k}x_ix_k
= \Bigl(\sum x_i\Bigr)^2 - 3\sum_{i<k}x_ix_k = -3p,
\]
using $j + j^2 = -1$ and $\sum_{i<k}x_ix_k = p$. Then
\[
u^3 + v^3 = (u+v)^3 - 3uv(u+v) = 27x_1^3 + 9p\cdot 3x_1
= 27\,(x_1^3 + px_1) = -27q .
\]
So $u^3, v^3$ are the roots of $Y^2 + 27qY - 27p^3 = 0$
(product $(uv)^3 = -27p^3$): $u^3 = \frac{-27q +
\sqrt{729q^2 + 108p^3}}2$, and $x_1 = \frac{u + v}3$ with $v =
-3p/u$: Cardano. Solvability of $S_3$ is the skeleton: the
tower $\Q \subseteq \Q(\delta) \subseteq \Q(\delta, j, u)$
adjoins first a square root ($\delta$, fixed field of $A_3$:
the step $S_3 \to S_3/A_3$), then a cube root ($u$, since $u^3
\in \Q(\delta, j)$: the step $A_3 \to \{e\}$) --- the derived
series $S_3 \supset A_3 \supset \{e\}$ made flesh.
\end{solution}

\begin{solution}{exo:b3:galois:9}
$\eta_0 + \eta_1 = \zeta_5 + \zeta_5^2 + \zeta_5^3 + \zeta_5^4 =
-1$ (sum of all $5$-th roots of unity is $0$). $\eta_0\eta_1 =
(\zeta + \zeta^4)(\zeta^2 + \zeta^3) = \zeta^3 + \zeta^4 +
\zeta^6 + \zeta^7 = \zeta^3 + \zeta^4 + \zeta + \zeta^2 = -1$
(indices mod $5$). So $\eta_0, \eta_1$ are the roots of $Y^2 + Y
- 1$: $\frac{-1 \pm \sqrt5}2$. Since $\eta_0 =
2\cos\frac{2\pi}5 > 0$: $\eta_0 = \frac{\sqrt5 - 1}2$, whence
$\cos\frac{2\pi}5 = \frac{\sqrt5 - 1}4$, and $\eta_1 =
\frac{-1-\sqrt5}2$. The group
$\operatorname{Gal}(\Q(\zeta_5)/\Q) \cong (\Z/5\Z)^\times$ is
cyclic of order $4$: it has a \emph{unique} subgroup of order
$2$ ($\{\pm 1\}$, i.e.\ $\zeta \mapsto \zeta^{\pm1}$), hence
$\Q(\zeta_5)$ has a unique quadratic subfield
(\cref{thm:b3:galois:fundamental}), which contains $\eta_0
\notin \Q$: it is $\Q(\eta_0) = \Q(\sqrt5)$. Constructibility:
$\cos\frac{2\pi}5$ lies in the quadratic tower $\Q \subseteq
\Q(\sqrt5)$, and $\zeta_5$ one quadratic step above:
\cref{thm:b3:galois:wantzel} constructs the pentagon.
\end{solution}

\begin{solution}{exo:b3:galois:10}
(a) Write $s = S^{1/p}$, $t = T^{1/p}$ (elements of a chosen
algebraic closure with $s^p = S$, $t^p = T$). $X^p - S$ is
irreducible over $K = \mathbb F_p(S, T) = (\mathbb
F_p(T))(S)$-fractions: Eisenstein at the prime element $S$ of
the UFD $\mathbb F_p(T)[S]$ (\cref{thm:b3:rings:criteria}). So
$[K(s):K] = p$; likewise $X^p - T$ is Eisenstein at $T$ over
$K(s) = \mathbb F_p(s)(T)$-fractions --- $T$ remains prime in
$\mathbb F_p(s)[T]$ --- giving $[L : K(s)] = p$ and $[L:K] =
p^2$. For $\alpha \in L$: $L = K[s, t]$, so $\alpha = \sum
c_{ij}s^it^j$ ($c_{ij} \in K$), and by the Frobenius morphism
$\alpha^p = \sum c_{ij}^p S^iT^j \in K$.

(b) If $L = K(\alpha)$, then $[K(\alpha):K] = p^2$; but
$\alpha^p = a \in K$ means $\alpha$ annihilates $X^p - a$, so
$\deg\pi_\alpha \leq p < p^2$: contradiction. No primitive
element: \cref{thm:b3:galois:primitive} genuinely needs
separability (here every $\pi_\alpha$ divides some $X^p - a =
(X - \alpha)^p$: purely inseparable).
\end{solution}

\begin{solution}{exo:b3:galois:11}
(a) Eisenstein at $2$ ($2 \mid 4, 2$; $4 \nmid 2$): $P$
irreducible. If $\alpha$ is a root, $[\Q(\alpha):\Q] = 5$
divides $[L:\Q] = \abs G$ ($L$ the splitting field): Cauchy
(\cref{thm:b3:groups:cauchy}) gives an element of order $5$ in
$G \leq S_5$; in $S_5$, only $5$-cycles have order $5$ (orders
are lcms of cycle lengths).

(b) $P'(x) = 5x^4 - 4$ vanishes at $\pm(4/5)^{1/4} \approx \pm
0.946$: one local maximum then one local minimum. Values: $P(-2)
= -22 < 0$, $P(0) = 2 > 0$, $P(1) = -1 < 0$, $P(2) = 26 > 0$:
three sign changes, and at most three real roots (two critical
points): exactly $3$ real roots, hence one pair of complex
conjugate roots. Take the splitting field $L$ inside $\C$:
complex conjugation maps $L$ to itself (it permutes the roots,
which generate $L$) and fixes $\Q$, so it defines an element of
$G$; it fixes the three real roots and swaps the other two: a
transposition.

(c) Let $\tau = (a\,b)$ and $\sigma$ a $5$-cycle in $G$. Some
power $\sigma^k$ sends $a$ to $b$ ($k \ne 0 \bmod 5$), and
$\sigma^k$ is again a $5$-cycle: renaming, assume $\sigma =
(1\,2\,3\,4\,5)$ and $\tau = (1\,2)$. Conjugating,
$\sigma^m\tau\sigma^{-m} = (\sigma^m(1)\ \sigma^m(2))$: the
adjacent transpositions $(1\,2), (2\,3), (3\,4), (4\,5),
(5\,1)$ all lie in $G$; adjacent transpositions generate $S_5$
(every transposition $(i\,j)$ is a product of adjacent ones,
and transpositions generate). So $G = S_5$, not solvable
(\cref{cor:b3:groups:snnotsolvable}), and
\cref{thm:b3:galois:solvable} concludes: $X^5 - 4X + 2$ is not
solvable by radicals.
\end{solution}

\begin{solution}{pb:b3:galois:1}
\textbf{1.} $\Phi_{17}$ is irreducible
(\cref{thm:b3:galois:cyclotomicirred}, or
\cref{ex:b3:rings:eisenstein} for prime index): $[L:\Q] =
\varphi(17) = 16$ and $G \cong (\Z/17\Z)^\times$, cyclic of
order $16$ (\cref{thm:b3:galois:cyclic}). Powers of $3$ mod
$17$:
\[
3,\ 9,\ 10,\ 13,\ 5,\ 15,\ 11,\ 16,\ 14,\ 8,\ 7,\ 4,\ 12,\ 2,\
6,\ 1
\]
--- sixteen distinct values: $3$ generates.

\textbf{2.} $G = \langle\sigma\rangle$ cyclic of order $16$;
$H_k = \langle\sigma^{2^k}\rangle$ has order $2^{4-k}$, and
$[H_k : H_{k+1}] = 2$. By the fundamental theorem
(\cref{thm:b3:galois:fundamental}), $L_k = L^{H_k}$ satisfy
$[L_k : \Q] = [G : H_k] = 2^k$: each $[L_{k+1}:L_k] = 2$.

\textbf{3.} $\zeta \in L = L_4$ sits atop a tower of quadratic
extensions of $\Q$: by \cref{thm:b3:galois:wantzel}, $\zeta$ is
constructible; the $17$-gon has vertices $\zeta^k$.

\textbf{4.} $\sigma^2$ multiplies exponents by $9$; the
exponents of $\eta_0$ are the even powers of $3$,
\[
\{3^{2k} \bmod 17\} = \{1, 9, 13, 15, 16, 8, 4, 2\},
\]
a set stable under multiplication by $9 = 3^2$; so $\eta_0$
(and likewise $\eta_1$) is fixed by $H_1 =
\langle\sigma^2\rangle$: $\eta_0, \eta_1 \in L_1$, a quadratic
field. $\sigma$ maps even powers to odd: it swaps $\eta_0,
\eta_1$. Hence $\eta_0 + \eta_1$ and $\eta_0\eta_1$ are fixed by
all of $G$: rational; $\eta_0, \eta_1$ are the roots of a
rational quadratic.

\textbf{5.} $\eta_0 + \eta_1 = \sum_{c=1}^{16}\zeta^c = -1$. The
product expands into $64$ terms $\zeta^{a + b}$, $a$ in the even
set, $b$ in the odd set. No term is $\zeta^0$: $b = -a$ is
impossible, because $-1 = 16 = 3^8$ is an \emph{even} power, so
$-a$ stays in the even set. Thus $\eta_0\eta_1 = \sum_{c \neq 0}
n_c\zeta^c$ with $\sum n_c = 64$; applying $\sigma$ fixes
$\eta_0\eta_1$ (it swaps the factors) and permutes the $\zeta^c$
transitively over all $c \neq 0$, so all $n_c$ are equal: $n_c
= 4$ and $\eta_0\eta_1 = 4\sum_{c\neq0}\zeta^c = -4$.

\textbf{6.} $\eta_{0,1}$ solve $Y^2 + Y - 4 = 0$: $\frac{-1 \pm
\sqrt{17}}2$. Numerically, pairing conjugate exponents,
$\eta_0 = 2\bigl(\cos\tfrac{2\pi}{17} + \cos\tfrac{4\pi}{17} +
\cos\tfrac{8\pi}{17} + \cos\tfrac{16\pi}{17}\bigr) \approx 1.56
> 0$: $\eta_0 = \frac{-1+\sqrt{17}}2$, $\eta_1 =
\frac{-1-\sqrt{17}}2$, and $L_1 = \Q(\eta_0) = \Q(\sqrt{17})$.

\textbf{7.} The exponent sets: $\beta_0$: $\{1, 13, 16, 4\}$ =
powers $3^{4k}$; $\beta_2$: $\{9, 15, 8, 2\}$ = $9 \times$ that
set. Union: the even set: $\beta_0 + \beta_2 = \eta_0$; likewise
$\beta_1 + \beta_3 = \eta_1$. Multiplication by $13 = 3^4$
stabilizes each $\beta_i$'s exponent set: fixed by $H_2 =
\langle\sigma^4\rangle$; and $\sigma^2$ ($\times 9$) sends
$\{1,13,16,4\}$ to $\{9, 15, 8, 2\}$: swaps $\beta_0,
\beta_2$.

\textbf{8.} Expanding $\beta_0\beta_2$, the sixteen exponent
sums
\[
\{1,13,16,4\} + \{9,15,8,2\} =
\{10,16,9,3,\ 5,11,4,15,\ 8,14,7,1,\ 13,2,12,6\}
\]
cover $1, \dots, 16$ exactly once: $\beta_0\beta_2 =
\sum_{c\ne0}\zeta^c = -1$. Applying $\sigma$ (which maps
$\beta_0 \mapsto \beta_1$, $\beta_2 \mapsto \beta_3$:
exponents $\times 3$): $\beta_1\beta_3 = \sigma(\beta_0\beta_2)
= -1$.

\textbf{9.} $\beta_0, \beta_2$ solve $Y^2 - \eta_0 Y - 1 = 0$,
so $\beta_0 = \frac{\eta_0 + \sqrt{\eta_0^2 + 4}}2$ (numerically
$\beta_0 = 2\cos\frac{2\pi}{17} + 2\cos\frac{8\pi}{17} \approx
2.05 > 0$, the $+$ sign). Likewise $\beta_1 = \frac{\eta_1 +
\sqrt{\eta_1^2 + 4}}2 \approx 0.344$ (numerical check fixes the
sign again). $L_2 = L_1(\beta_0)$, quadratic over $L_1$.

\textbf{10.} $\gamma_0 + \gamma_1 = \zeta + \zeta^{16} +
\zeta^{13} + \zeta^4 = \beta_0$. And
\[
\gamma_0\gamma_1 = (\zeta + \zeta^{16})(\zeta^{13} + \zeta^4)
= \zeta^{14} + \zeta^{5} + \zeta^{12} + \zeta^{3} = \beta_1 .
\]
So $\gamma_0, \gamma_1$ solve $Y^2 - \beta_0Y + \beta_1 = 0$;
numerically $\gamma_0 = 2\cos\frac{2\pi}{17} \approx 1.865 >
\gamma_1 \approx 0.185$: $\gamma_0 = \frac{\beta_0 +
\sqrt{\beta_0^2 - 4\beta_1}}2$.

\textbf{11.} Chaining:
\[
\eta_0 = \frac{-1 + \sqrt{17}}2,
\quad
\beta_0 = \frac{\eta_0 + \sqrt{\eta_0^2 + 4}}2,
\quad
\beta_1 = \frac{\eta_1 + \sqrt{\eta_1^2 + 4}}2,
\quad
\cos\frac{2\pi}{17} = \frac{\beta_0 + \sqrt{\beta_0^2 -
4\beta_1}}4 .
\]
Numerically: $\sqrt{17} \approx 4.1231$, $\eta_0 \approx
1.5616$, $\eta_1 \approx -2.5616$, $\beta_0 \approx 2.0494$,
$\beta_1 \approx 0.3441$, $\beta_0^2 - 4\beta_1 \approx
2.8234$, and $\cos\frac{2\pi}{17} \approx \frac{2.0494 +
1.6803}4 \approx 0.93242$ --- against $\cos\frac{2\pi}{17} =
0.93247\dots$: the small discrepancy is rounding in the
intermediate displays; carrying more digits reproduces
$0.932472$.

\textbf{12.} The construction needed $[L:\Q] = p - 1$ to be a
power of $2$, so that a full chain of index-$2$ subgroups
exists. If $p = 2^m + 1$ is prime and $m = ab$ with $a$ odd $>
1$: $x + 1 \mid x^a + 1$ at $x = 2^b$ shows $2^b + 1$ properly
divides $p$ --- impossible. So $m$ is a power of $2$: $p =
2^{2^t} + 1$, a \emph{Fermat prime} ($3, 5, 17, 257, 65537$,
\dots). Conversely for such $p$, $\varphi(p) = 2^{2^t}$ and the
argument of questions 1--3 (or
\cref{cor:b3:galois:impossible}) applies: the regular $p$-gon
is constructible iff $p$ is a Fermat prime.

\textbf{13.} $\varphi(n)$ is a power of $2$ exactly when $n =
2^a p_1\cdots p_r$ with distinct Fermat primes $p_i$
(multiplicativity of $\varphi$; an odd prime power $p^k$, $k
\geq 2$, contributes the factor $p \nmid 2^m$). For $n \leq
20$, the constructible regular $n$-gons are
\[
n = 3, 4, 5, 6, 8, 10, 12, 15, 16, 17, 20
\]
with respective values
\[
\varphi(n) = 2,\ 2,\ 4,\ 2,\ 4,\ 4,\ 4,\ 8,\ 8,\ 16,\ 8 .
\]
The impossible ones are $n = 7, 9, 11, 13, 14, 18, 19$, where
$\varphi(n) = 6, 6, 10, 12, 6, 6, 18$ has an odd prime factor.
\end{solution}
